\begin{problem}[第32届全国中学生物理竞赛决赛]{50}
俄国火箭专家齐奥尔科夫斯基将火箭发射过程进行模型简化，得出了最早的理想火箭方程，为近代火箭、导弹工程提供了理论依据。该简化模型为：待发射火箭静止于惯性参考系S中的某点，忽略火箭所受的地球引力等外力的作用，火箭（包含燃料）的初始静止质量为 $M_{i}$；在t=0时刻点火，火箭向左排出气体，气体相对于火箭的速度恒为 $V_{e}$，使得火箭向右发射。在S系中观测，火箭的速度为 $V_{1}$，排出气体的速度为 $V_{2}$，如图所示。根据此模型，火箭运行一段时间后，当其静止质量由 $M_{i}$ 变为M时，
\begin{subproblem}
用牛顿力学推导火箭所获得的速度 $V_{1}$ 与质量比 $M/M_{i}$ 之间的关系；
\end{subproblem}
\begin{subproblem}
用狭义相对论力学推导火箭所获得的速度 $V_{1}$ 与质量比 $M/M_{i}$ 之间的关系；
\end{subproblem}
\begin{subproblem}
当 $V_{1}$ 远小于真空中的光速c时，计算以上两种结果之差，保留至 $(V_{1}/c)^{2}$ 项。
\end{subproblem}
\end{problem}